%%
% BIThesis 读书报告模板 —— 使用 XeLaTeX 编译 The BIThesis Template for Reading Report
% This file has no copyright assigned and is placed in the Public Domain.
% Compile with: xelatex -> biber -> xelatex -> xelatex
%%

% 请勿删除下面两行注释，以免影响编译。
% !TeX program = xelatex
% !BIB program = biber


\documentclass[type=bachelor]{bithesis}

% 此处仅列出常用的配置。全部配置用法请见「bithesis.pdf」手册。
\BITSetup{
  cover = {
    % 封面需要「北京理工大学」字样图片，如无必要请勿修改该项。
    headerImage = images/header.png,
    % 封面标题需要“华文细黑”，如无必要请勿修改该项。
    xiheiFont = STXIHEI.TTF,
    % 可用以下参数自定义封面日期。
    % date = 1955年11月,
  },
  info = {
    % 想要删除某项封面信息，直接删除该项即可。
    % 想要让某项封面信息留空（但是保留下划线），请传入空白符组成的字符串，如"{~}"。
    % 如需要换行，则用 “\\” 符号分割。
    title = 北京理工大学本科生读书报告题目,
    semester = {1840-1841学年第一学期},
    school = 信息与电子学院,
    author = 罗辑,
    studentId = 11xxxxxxxx,
    course = 宇宙社会学,
    teacher = 叶文洁,
    % 以下用于隐藏字段
    titleEn = {},
  },
  style = {
    head = 北京理工大学本科生读书报告,
    headline = 本科生读书报告,
    % 开启 Windows 平台下的中易宋体伪粗体。
    % windowsSimSunFakeBold = true,
    %
    % 开启该选项后，将用 Times New Roman 的开源字体 TeX Gyre Termes 作为正文字体。
    % 这个选项适用于以下情况：
    % 1. 不想在系统中安装 Times New Roman。
    % 2. 在 Linux/macOS 下遇到 `\textsc` 无法正常显示的问题。
    % betterTimesNewRoman = true,
  },
  misc = {
    % 微调表格行间距
    tabularRowSeparation = 1.25,
  },
}

% 这份模板提供了代码块示例，采用 listings 宏包及其预定义样式。
% 使用代码块时，也可选用自己的喜欢的其他宏包，如 minted：https://bithesis.bitnp.net/faq/minted.html
% 如果不需要，直接删除即可。
\usepackage{listings}


% 大部分关于参考文献样式的修改，都可以通过此处的选项进行配置。
% 详情请搜索「biblatex-gb7714-2015 文档」进行阅读。
\usepackage[
  backend=biber,
  style=gb7714-2015,
  gbalign=gb7714-2015,
  gbnamefmt=lowercase,
  gbpub=false,
  doi=false,
  url=false,
  eprint=false,
  isbn=false,
]{biblatex}

% 参考文献引用文件位于 misc/ref.bib
\addbibresource{misc/ref.bib}

% 文档开始
\begin{document}

% 标题页面：如无特殊需要，本部分无需改动
\MakeCover

% 前置内容定义
\frontmatter

\MakeTOC

% 正文开始
\mainmatter

% 在这里引用各章 TeX 文件，按需添加
%%
% BIThesis 读书报告模板 —— 使用 XeLaTeX 编译 The BIThesis Template for Reading Report
% This file has no copyright assigned and is placed in the Public Domain.
%%

% 第一章节

\chapter{一级题目}

\section{二级题目}
% 这里插入一个参考文献，仅作参考

\subsection{三级题目}

正文……\parencite{yuFeiJiZongTiDuoXueKeSheJiYouHuaDeXianZhuangYuFaZhanFangXiang2008}……\cite{Hajela2012Application}

\textcolor{blue}{正文部分：宋体、小四；正文行距：22磅；间距段前段后均为0行。阅后删除此段。}

\textcolor{blue}{图、表居中，图注标在图下方，表头标在表上方，宋体、五号、居中，1.25倍行距，间距段前段后均为0行，图表与上下文之间各空一行。阅后删除此段。}

\textcolor{blue}{\underline{\underline{图-示例：（阅后删除此段）}}}


\begin{figure}[htbp]
  \centering
  \includegraphics[]{images/bit_logo.png}
  \caption{标题序号}\label{标题序号} % label 用来在文中索引
\end{figure}

\textcolor{blue}{\underline{\underline{表-示例：（阅后删除此段）}}}
% 三线表
\begin{table}[htbp]
  \centering
  \caption{统计表}\label{统计表}
  \begin{tabular}{*{5}{>{\centering\arraybackslash}p{2cm}}} \toprule
    项目    & 产量    & 销量    & 产值   & 比重    \\ \midrule
    手机    & 1000  & 10000 & 500  & 50\%  \\
    计算机   & 5500  & 5000  & 220  & 22\%  \\
    笔记本电脑 & 1100  & 1000  & 280  & 28\%  \\ \midrule
    合计    & 17600 & 16000 & 1000 & 100\% \\ \bottomrule
    \end{tabular}
\end{table}

\textcolor{blue}{公式标注应于该公式所在行的最右侧。对于较长的公式只可在符号处（+、-、*、/、$\leqslant$ $\geqslant$ 等）转行。在文中引用公式时，在标号前加“式”，如式（1-2）。阅后删除此
段。}

\textcolor{blue}{公式-示例：（阅后删除此段）}
% 公式上下不要空行，置于同一个段落下即可，否则上下距离会出现高度不一致的问题
\begin{equation}
    LRI=1\ ∕\ \sqrt{1+{\left(\frac{{\mu }_{R}}{{\mu }_{s}}\right)}^{2}{\left(\frac{{\delta }_{R}}{{\delta }_{s}}\right)}^{2}}
\end{equation}


\section{字体效果表格}

% 列：Regular、Italic、Bold、Bold Italic
% 行：宋体、黑体、楷体、Serif、Sans Serif、Typewriter、Math

\begin{table}[htb]
    \centering
    \caption{字体效果表格}
    \begin{tabular}{@{}lllll@{}}
    \toprule
               & Regular & Bold & Italic & Bold Italic \\ \midrule
      宋体       & 宋体      & \colorbox{orange}{\textbf{宋体粗体}} & \textit{楷体}     &   \colorbox{gray}{\textbf{\textit{楷书粗斜体}}}  \\
      黑体         & {\heiti{}黑体}      & \textbf{\heiti{}黑体粗体} &   \textit{\heiti{}黑体斜体}     & \colorbox{gray}{\textit{\textbf{\heiti{}黑体粗斜体}}}   \\
      楷体         & {\kaishu{}楷书}      & \textbf{\kaishu{}楷书粗体} & \textit{\kaishu{}斜体楷体} &  \colorbox{gray}{\textbf{\textit{\kaishu{}楷书粗斜体}}}    \\
    Serif(Roman/Normal)      &    Regular    &  \textbf{Bold}  &    \textit{Italic}    &     \textbf{\textit{Bold Italic}}    \\
    Sans Serif &  \textsf{Regular}       &  \textbf{\textsf{Bold}}    &  \textit{\textsf{Bold}}   &    \textbf{\textit{\textsf{Bold}}}   \\
    % 有些字体缺少特定变体，LaTeX会抛出警告，同时尽量替换为相近变体。
    % 若编译结果能接受，可忽略之。
    % 例如此处Typewriter代表的Latin Modern Sans Typewriter（lmtt）字体没有bold extended italic（bx/it）变体，
    % 所以`\textbf{\textit{\texttt{…}}}`会被替换为bold slant（b/sl）变体，并引发如下警告。（其中TU指字体编码）
    % Font shape `TU/lmtt/bx/it' in size <…> not available. Font shape `TU/lmtt/b/sl' tried instead.
    Typewriter &  \texttt{Regular}       &  \textbf{\texttt{Bold}}    &  \textit{\texttt{Bold}}   &    \textbf{\textit{\texttt{Bold}}}   \\
    Math       &   $\mathnormal{Regular} \mathrm{Roman}$  & $\mathbf{Bold}$   &    $\mathit{Italic}$    &  $\mathbf{\mathit{Bold Italic}}$    \\ \bottomrule
    \end{tabular}
\end{table}

\begin{itemize}[nosep]
  \item \colorbox{orange}{宋体粗体}在 Windows 下会成为黑体。这是因为 Windows 的中易宋体没有粗体字重而进行的妥协。
    如果想要获得宋体粗体的样式，请在配置中开启伪粗体选项。
  \item \colorbox{gray}{粗斜体}的效果是因操作系统字体而定的，中文写作中不会使用这种字形，可以忽略。
\end{itemize}


\section{常见问题和疑难解答}

如果您遇到\href{https://bithesis.bitnp.net/faq/char-missing.html}{生僻字无法显示}、
\href{https://bithesis.bitnp.net/faq/enumitem-nosep.html}{列表项间距过大}、
\href{https://bithesis.bitnp.net/faq/longtable.html}{三线表需要跨页}等问题，
请参考\textcolor{magenta}{\href{https://bithesis.bitnp.net/faq/}{在线文档的「疑难杂症」部分}}。

%%
% BIThesis 读书报告模板 —— 使用 XeLaTeX 编译 The BIThesis Template for Reading Report
% This file has no copyright assigned and is placed in the Public Domain.
%%

\chapter{另一个章节}

\section{代码片段}

\begin{lstlisting}[language=Python, caption={Python Code}, label={lst:pythonfile}]
import numpy as np

def incmatrix(genl1,genl2):
    m = len(genl1)
    n = len(genl2)
    M = None #to become the incidence matrix
    VT = np.zeros((n*m,1), int)  #dummy variable

    #compute the bitwise xor matrix
    M1 = bitxormatrix(genl1)
    M2 = np.triu(bitxormatrix(genl2),1)

    for i in range(m-1):
        for j in range(i+1, m):
            [r,c] = np.where(M2 == M1[i,j])
            for k in range(len(r)):
                VT[(i)*n + r[k]] = 1;
                VT[(i)*n + c[k]] = 1;
                VT[(j)*n + r[k]] = 1;
                VT[(j)*n + c[k]] = 1;

                if M is None:
                    M = np.copy(VT)
                else:
                    M = np.concatenate((M, VT), 1)

                VT = np.zeros((n*m,1), int)

    return M
\end{lstlisting}


% 后置内容
\backmatter

% 结论：在相应 TeX 文件撰写
%%
% BIThesis 读书报告模板 —— 使用 XeLaTeX 编译 The BIThesis Template for Reading Report
% This file has no copyright assigned and is placed in the Public Domain.
%% 

\begin{conclusion}
  % 结论部分尽量不使用 \subsection 二级标题，只使用 \section 一级标题

  % 这里插入一个参考文献，仅作参考
  本文结论……\cite{张伯伟2002全唐五代诗格会考}。

  \textcolor{blue}{结论作为正文的最后部分单独排写，但不加章号。阅后删除此段。}

  \textcolor{blue}{结论正文样式与文章正文相同：宋体、小四；行距：22磅；间距段前段后均为0行。阅后删除此段。}
\end{conclusion}

% 参考文献：
% 添加文献时，请按 BibTeX 格式添加至 misc/ref.bib，并在正文所需位置使用 \cite{…} 引用。
% 如无特殊需要，无需改动相应 TeX 文件。
%%
% BIThesis 读书报告模板 —— 使用 XeLaTeX 编译 The BIThesis Template for Reading Report
% This file has no copyright assigned and is placed in the Public Domain.
%%

\begin{bibprint}

% \printbibliography[heading=none]
% 正式使用时，请启用上方语句以输出所有的参考文献，并删除/注释下方示例内容。
% 删除后仍可参考 https://bithesis.bitnp.net/faq/bib-entry.html 的在线版本。

% -------------------------------- 示例内容（正式使用时请删除） ------------------------------------- %

% 抑制多次调用 \printbibliography 的 warning，只有示例代码会需要此语句。
\BiblatexSplitbibDefernumbersWarningOff

\textcolor{blue}{参考文献书写规范}

\textcolor{blue}{参考国家标准《信息与文献参考文献著录规则》【GB/T 7714—2015】，参考文献书写规范如下：}

\textcolor{blue}{\textbf{1. 文献类型和标识代码}}

\textcolor{blue}{普通图书：M}\qquad\textcolor{blue}{会议录：C}\qquad\textcolor{blue}{汇编：G}\qquad\textcolor{blue}{报纸：N}

\textcolor{blue}{期刊：J}\qquad\textcolor{blue}{学位论文：D}\qquad\textcolor{blue}{报告：R}\qquad\textcolor{blue}{标准：S}

\textcolor{blue}{专利：P}\qquad\textcolor{blue}{数据库：DB}\qquad\textcolor{blue}{计算机程序：CP}\qquad\textcolor{blue}{电子公告：EB}

\textcolor{blue}{档案：A}\qquad\textcolor{blue}{舆图：CM}\qquad\textcolor{blue}{数据集：DS}\qquad\textcolor{blue}{其他：Z}

\textcolor{blue}{\textbf{2. 不同类别文献书写规范要求}}

\textcolor{blue}{\textbf{期刊}}

\noindent\textcolor{blue}{[序号] 主要责任者. 文献题名[J]. 刊名, 出版年份, 卷号(期号): 起止页码. }
\cite{yuFeiJiZongTiDuoXueKeSheJiYouHuaDeXianZhuangYuFaZhanFangXiang2008, Hajela2012Application}

\printbibliography [type=article,heading=none]

\textcolor{blue}{\textbf{普通图书}}

\noindent\textcolor{blue}{[序号] 主要责任者. 文献题名[M]. 出版地: 出版者, 出版年: 起止页码. }
\cite{张伯伟2002全唐五代诗格会考, OBRIEN1994Aircraft}

\printbibliography [keyword={book},heading=none]

\textcolor{blue}{\textbf{会议论文集}}

\noindent\textcolor{blue}{[序号] 主要责任者．题名:其他题名信息[C]. 出版地: 出版者, 出版年. }
\cite{雷光春2012}

\printbibliography [type=proceedings,heading=none]

\textcolor{blue}{\textbf{专著中析出的文献}}

\noindent\textcolor{blue}{[序号] 析出文献主要责任者. 析出题名[M]//专著主要责任者. 专著题名. 出版地: 出版者, 出版年: 起止页码. }
\cite{白书农}

\printbibliography [type=inbook,heading=none]

\textcolor{blue}{\textbf{学位论文}}

\noindent\textcolor{blue}{[序号] 主要责任者. 文献题名[D]. 保存地: 保存单位, 年份. }
\cite{zhanghesheng, Sobieski}

\printbibliography [keyword={thesis},heading=none]

\textcolor{blue}{\textbf{报告}}

\noindent\textcolor{blue}{[序号] 主要责任者. 文献题名[R]. 报告地: 报告会主办单位, 年份. }
\cite{fengxiqiao, Sobieszczanski}

\printbibliography [keyword={techreport},heading=none]

\textcolor{blue}{\textbf{专利文献}}

\noindent\textcolor{blue}{[序号] 专利所有者. 专利题名:专利号[P]. 公告日期或公开日期[引用日期]. 获取和访问路径. 数字对象唯一标识符.}
\cite{jiangxizhou}

\printbibliography [type=patent,heading=none]

\textcolor{blue}{\textbf{国际、国家标准}}

\noindent\textcolor{blue}{[序号] 主要责任人. 题名: 其他题名信息[S]. 出版地: 出版者, 出版年: 引文页码.}
\cite{GB/T3792.4-2009}

\printbibliography [keyword={standard},heading=none]

\textcolor{blue}{\textbf{报纸文章}}

\noindent\textcolor{blue}{[序号] 主要责任者. 文献题名[N]. 报纸名, 年(期): 页码. }
\cite{xiexide}

\printbibliography [keyword={newspaper},heading=none]

\textcolor{blue}{\textbf{电子文献}}

\noindent\textcolor{blue}{[序号] 主要责任者. 电子文献题名[文献类型/载体类型]. (发表或更新日期) [引用日期]. 获取和访问路径. 数字对象唯一标识符. }
\cite{yaoboyuan}

\printbibliography [keyword={online},heading=none]

\textcolor{blue}{关于参考文献的未尽事项可参考国家标准《信息与文献参考文献著录规则》（GB/T 7714—2015）}

\end{bibprint}


\end{document}
